\section{Experiments}

\subsection{Datasets and Training Setup}

We evaluate on the 200-video SHOT test split \cite{zhu2023autoshot}, all 11
BBC Planet Earth videos \cite{baraldi2015shot}, and the 500-video ClipShots
test split \cite{tang2018clipshots}. SHOT is the primary short-video
benchmark; BBC tests long documentary footage with many smooth gradual
transitions; ClipShots tests open-world transfer. Table~\ref{tab:datasets}
summarizes their characteristics.

\begin{table}[!t]
\caption{Dataset Characteristics Used for SBD Evaluation}
\label{tab:datasets}
\centering
\scriptsize
\renewcommand{\arraystretch}{1.08}
\setlength{\tabcolsep}{3pt}
\begin{tabular}{L{0.95in}C{0.78in}C{0.85in}C{0.62in}}
\toprule
\textbf{Property} & \textbf{BBC} & \textbf{ClipShots} & \textbf{SHOT} \\
\midrule
Content type & Documentary & Open-world video & Short video \\
Total videos & 11 & 4,039 & 853 \\
Evaluation videos & 11 & 500 & 200 \\
Avg. video length & 2,945 s & 237 s & 39.5 s \\
Avg. shot length & 6.57 s & 15.34 s & 2.59 s \\
Hard cuts & Many & 128,000+ & 8,890 \\
Gradual transitions & Many & 38,000+ & 2,716 \\
Main challenge & Gradual trans. & Diverse effects & Fast editing \\
\bottomrule
\end{tabular}
\end{table}

AutoShotV2 trains only the classification head with Adam at a constant
learning rate; Table~\ref{tab:hyperparams} lists the hyperparameters and how
each was selected. Performance is measured by precision, recall, and F1,
micro-averaged from transition-level TP/FP/FN counts with a matching
tolerance of $\pm\PaperMatchTolerance$ frames.

\begin{table}[!t]
\caption{AutoShotV2 Hyperparameters}
\label{tab:hyperparams}
\centering
\small
\renewcommand{\arraystretch}{1.08}
\begin{tabular}{L{1.25in}C{0.7in}L{0.85in}}
\toprule
\textbf{Component} & \textbf{Value} & \textbf{Selection} \\
\midrule
Focal Loss $\gamma$ & 1.0 & Grid search \\
Focal Loss $\alpha$ & 0.5 & Grid search \\
Learning rate & $10^{-5}$ & Adam, constant \\
Weight decay & $10^{-4}$ & Fixed \\
Batch size & 512 & Frame-level \\
Maximum epochs & 20 & Fixed budget \\
Dropout & 0.5 & Fixed \\
Many-hot weight & 0.3 & Validation \\
Gaussian $\sigma$ & \PaperGaussianSigma{} & Validation \\
Deploy temperature $T^{*}$ & \PaperDeployTemperature{} & Validation BCE \\
\bottomrule
\end{tabular}
\end{table}

\subsection{Evaluation Tracks and Reproducibility Tiers}
\label{sec:tracks}

The paper reports two evaluation tracks with different reproducibility
tiers. \emph{Track 1 (deployed checkpoint)} provides the headline numbers:
the head trained for deployment with the fixed configuration
$T^{*}=\PaperDeployTemperature$, $\sigma=\PaperGaussianSigma$,
$\theta=\PaperDeployThreshold$. Its per-dataset result JSON files are
tracked and checksummed, but its test logits were not retained, so
logit-level analyses cannot be re-run for this checkpoint. \emph{Track 2
(controlled replication)} retrains the head under the frozen validation-only
protocol of Section~\ref{sec:protocol} (configuration B4) with full logit
caches; every confidence interval, paired bootstrap, calibration diagnostic,
ablation, and seed study in this paper is anchored on this track. The two
tracks use different validation pools, so their temperatures and absolute
scores differ slightly; reporting both keeps the headline faithful to the
deployed system while keeping every statistical claim verifiable from
tracked artifacts.

\subsection{Main Results}

Table~\ref{tab:overall} reports F1 across SHOT, BBC, and ClipShots. With the
fixed deployment configuration, AutoShotV2 reaches
\PaperDeployShotFOne{} on SHOT and \PaperDeployBBCFOne{} on BBC, improving
over the reproduced AutoShot baseline by \PaperDeployVsAutoShotPP{} F1
points and over the reported TransNetV2 baseline by
\PaperDeployVsTransNetPP{} points on SHOT (Fig.~\ref{fig:bars}). On
ClipShots the fixed configuration reaches \PaperDeployClipFOne{}; a separate
test-set threshold sweep reaches \PaperDeployClipBestFOne{} at
$\theta=0.15$, reported for transparency but not treated as the fixed
deployment result. Reported baselines use their authors' protocols, whereas
reproduced baselines use our evaluation pipeline; the TransNetV2
reproduction was not evaluated on SHOT.

\begin{table}[!t]
\caption{F1-Score Comparison Across Evaluation Datasets}
\label{tab:overall}
\centering
\small
\renewcommand{\arraystretch}{1.08}
\resizebox{\columnwidth}{!}{%
\begin{tabular}{L{2.15in}C{0.55in}C{0.55in}C{0.7in}}
\toprule
\textbf{Method} & \textbf{SHOT} & \textbf{BBC} & \textbf{ClipShots} \\
\midrule
\PaperMainResultRows
\bottomrule
\end{tabular}
}
\end{table}

\begin{figure}[!t]
\centering
\begin{tikzpicture}[font=\small]
\pgfmathsetmacro{\vone}{(\PaperDeployShotFOne-0.70)*20}
\pgfmathsetmacro{\vtwo}{(\PaperAutoShotShotFOne-0.70)*20}
\pgfmathsetmacro{\vthree}{(\PaperTransNetShotFOne-0.70)*20}
\fill[red!70] (0,2.0) rectangle (\vone,2.5);
\node[right] at (\vone,2.25) {\PaperDeployShotFOne};
\node[left] at (0,2.25) {AutoShotV2};
\fill[green!60!black] (0,1.1) rectangle (\vtwo,1.6);
\node[right] at (\vtwo,1.35) {\PaperAutoShotShotFOne};
\node[left] at (0,1.35) {AutoShot};
\fill[blue!60] (0,0.2) rectangle (\vthree,0.7);
\node[right] at (\vthree,0.45) {\PaperTransNetShotFOne};
\node[left] at (0,0.45) {TransNetV2};
\draw[-{Stealth[length=4pt]}] (0,0) -- (5.4,0) node[below left] {F1};
\foreach \x/\lab in {0/0.70, 1/0.75, 2/0.80, 3/0.85, 4/0.90, 5/0.95} {
  \draw (\x,0) -- (\x,-0.07) node[below, font=\scriptsize] {\lab};
}
\draw (0,0) -- (0,2.7);
\end{tikzpicture}
\caption{F1-score comparison on SHOT. AutoShotV2 improves the reproduced
AutoShot baseline by \PaperDeployVsAutoShotPP{} points and the reported
TransNetV2 baseline by \PaperDeployVsTransNetPP{} points.}
\label{fig:bars}
\end{figure}

\subsection{Controlled Component Analysis}

All remaining analyses use Track 2: the controlled B4 replication whose
protocol is summarized in Table~\ref{tab:protocol}.

\begin{table}[!t]
\caption{Fixed Reproducible Evaluation Protocol (Controlled Track)}
\label{tab:protocol}
\centering
\small
\renewcommand{\arraystretch}{1.08}
\begin{tabular}{L{1.45in}L{1.55in}}
\toprule
\textbf{Item} & \textbf{Setting} \\
\midrule
Training objective & BCE, one-hot labels \\
Training / data seed & 42 / 42 \\
Training budget & 20 epochs, batch size 512 \\
Sampling & At most 160 frames/video, 3 negatives/positive \\
Validation size & 200 videos \\
Validation selection & 5 source-stratified folds \\
Selection criterion & Mean macro-F1 across sources/folds \\
Selected temperature & \PaperBFourTemperature \\
Gaussian smoothing & $\sigma=\PaperGaussianSigma$ \\
Decision threshold & $\theta=\PaperDeployThreshold$ \\
Matching tolerance & $\pm\PaperMatchTolerance$ frames \\
Bootstrap & 10,000 video resamples, seed 42 \\
\bottomrule
\end{tabular}
\end{table}

Table~\ref{tab:main_results} compares configurations that share the controlled
Phase 2 protocol. B4 is selected exclusively by the frozen five-fold
validation manifest. Relative to A1, B4 improves SHOT by
\PaperBFourVsAOneShotPP{} F1 points and ClipShots by
\PaperBFourVsAOneClipPP{} points, while BBC is unchanged after rounding.
Temperature scaling alone preserves the validation-selected operating point;
Gaussian smoothing is the component that changes boundary decisions.

\begin{table}[!t]
\caption{Protocol-Matched F1 Scores (B4: Selected Configuration)}
\label{tab:main_results}
\centering
\small
\renewcommand{\arraystretch}{1.08}
\resizebox{\columnwidth}{!}{%
\begin{tabular}{L{2.15in}C{0.55in}C{0.55in}C{0.7in}}
\toprule
\textbf{Configuration} & \textbf{SHOT} & \textbf{BBC} & \textbf{ClipShots} \\
\midrule
\PaperProtocolMatchedRows
\bottomrule
\end{tabular}
}
\end{table}

Table~\ref{tab:confidence} reports the corresponding B4 operating points and
video-level bootstrap intervals. BBC has only 11 sampling units, so its
interval should be interpreted cautiously. ClipShots has high recall but lower
precision, indicating cross-domain false positives rather than a general
failure to detect transitions.

\begin{table}[!t]
\caption{B4 Operating Point and Video-Level Bootstrap 95\% Confidence Intervals}
\label{tab:confidence}
\centering
\scriptsize
\renewcommand{\arraystretch}{1.08}
\setlength{\tabcolsep}{2pt}
\begin{tabular}{L{0.48in}C{0.34in}C{0.95in}C{0.46in}C{0.46in}C{0.35in}}
\toprule
\textbf{Dataset} & \textbf{F1} & \textbf{F1 95\% CI} & \textbf{Precision} & \textbf{Recall} & \textbf{Videos} \\
\midrule
\PaperConfidenceRows
\bottomrule
\end{tabular}
\end{table}

Table~\ref{tab:paired_delta} directly bootstraps the per-video difference
between A1 at its validation-selected threshold and B4 at its frozen operating
point. The intervals support a positive B4 delta on SHOT and ClipShots. The BBC
interval contains zero, so we make no improvement claim for that dataset.

\begin{table}[!t]
\caption{Paired Video Bootstrap Difference, B4 Minus A1}
\label{tab:paired_delta}
\centering
\small
\renewcommand{\arraystretch}{1.08}
\resizebox{\columnwidth}{!}{%
\begin{tabular}{L{0.8in}C{0.6in}C{1.25in}C{0.65in}}
\toprule
\textbf{Dataset} & \textbf{$\Delta$F1} & \textbf{95\% CI} & \textbf{Excl. 0} \\
\midrule
\PaperPairedDeltaRows
\bottomrule
\end{tabular}
}
\end{table}

\subsection{Calibration Diagnostics}

We evaluate frame probabilities on the combined 200-video validation split
using negative log-likelihood (NLL), Brier score, expected calibration error
with 15 equal-width bins (ECE), equal-mass adaptive ECE (AECE), and
class-balanced ECE (B-ECE). Table~\ref{tab:calibration_metrics} shows that
temperature scaling improves all diagnostics. The much larger B-ECE reveals
that aggregate frame metrics are dominated by negative frames. Validation
values remain diagnostic because the same split fits temperature.

\begin{table}[!t]
\caption{Frame-Level Validation Calibration Diagnostics}
\label{tab:calibration_metrics}
\centering
\scriptsize
\renewcommand{\arraystretch}{1.08}
\setlength{\tabcolsep}{2pt}
\begin{tabular}{L{0.85in}C{0.45in}C{0.45in}C{0.45in}C{0.45in}C{0.45in}}
\toprule
\textbf{Variant} & \textbf{NLL} & \textbf{Brier} & \textbf{ECE} & \textbf{AECE} & \textbf{B-ECE} \\
\midrule
\PaperCalibrationMetricRows
\bottomrule
\end{tabular}
\end{table}

\begin{figure}[!t]
\centering
\includegraphics[width=\columnwidth]{reliability_diagram.png}
\caption{Reliability diagram on the combined validation split. Temperature
scaling reduces aggregate error. Equal-mass bins expose sparsely populated
confidence regions.}
\label{fig:reliability}
\end{figure}

Table~\ref{tab:test_calibration} reports untouched-test calibration using the
temperature selected on validation. ClipShots has the largest B-ECE despite a
small aggregate ECE, consistent with domain-shifted false positives.

\begin{table}[!t]
\caption{Test Calibration at the Validation-Selected Temperature}
\label{tab:test_calibration}
\centering
\scriptsize
\renewcommand{\arraystretch}{1.08}
\setlength{\tabcolsep}{2pt}
\begin{tabular}{L{0.45in}C{0.47in}C{0.47in}C{0.47in}C{0.47in}C{0.47in}}
\toprule
\textbf{Dataset} & \textbf{NLL} & \textbf{Brier} & \textbf{ECE} & \textbf{AECE} & \textbf{B-ECE} \\
\midrule
\PaperTestCalibrationRows
\bottomrule
\end{tabular}
\end{table}

\subsection{Controlled Training Ablation}
\label{sec:ablation}

Table~\ref{tab:ablation} reports the outcome of the design space of
Section~\ref{sec:objectives}. A2, A3, and B1 show that Focal Loss and
many-hot supervision do not improve the A1 control. B5 combines these
training changes with B4 post-processing and is numerically similar on SHOT
but slightly weaker on BBC and ClipShots. The gain therefore comes from
calibrated post-processing rather than the more complex training objective.
A0 is omitted because its per-dataset best-threshold protocol is not a
direct controlled comparison.

\begin{table}[!t]
\caption{Controlled Phase 2 Ablation (A0 Historical Baseline Excluded)}
\label{tab:ablation}
\centering
\scriptsize
\renewcommand{\arraystretch}{1.04}
\setlength{\tabcolsep}{2pt}
\begin{tabular}{L{1.7in}C{0.4in}C{0.4in}C{0.4in}}
\toprule
\textbf{Configuration} & \textbf{SHOT} & \textbf{BBC} & \textbf{ClipShots} \\
\midrule
\PaperControlledAblationRows
\bottomrule
\end{tabular}
\end{table}

\subsection{Efficiency and Reproducibility}

The backbone remains frozen, so only the Phase 2 head is optimized.
Table~\ref{tab:efficiency} reports parameter and run-artifact statistics.
The CPU post-processing measurement covers temperature scaling and Gaussian
smoothing on cached logits. End-to-end video decoding and backbone inference
are excluded. The original run artifact did not retain elapsed training time,
so we do not estimate or reconstruct it.

\begin{table}[!t]
\caption{Efficiency and Run-Artifact Statistics}
\label{tab:efficiency}
\centering
\small
\renewcommand{\arraystretch}{1.06}
\begin{tabular}{L{1.9in}L{1.05in}}
\toprule
\textbf{Item} & \textbf{Value} \\
\midrule
\PaperEfficiencyRows
\bottomrule
\end{tabular}
\end{table}

\subsection{ClipShots Transition-Type Recall}

We regenerate the ClipShots breakdown directly from the B4 logits and
\texttt{test.json}. One malformed annotation with reversed endpoints is
normalized by sorting the pair. A ground-truth interval is a cut when its
width is at most one frame and gradual otherwise. Predictions are
class-agnostic, so they are matched once and recall is grouped by the matched
ground-truth type. Type-specific precision and F1 are not reported because the
model does not predict a transition type.

\begin{table}[!t]
\caption{ClipShots Recall by Ground-Truth Transition Type}
\label{tab:clipshots_breakdown}
\centering
\small
\renewcommand{\arraystretch}{1.08}
\resizebox{\columnwidth}{!}{%
\begin{tabular}{L{0.8in}C{0.7in}C{0.65in}C{0.65in}C{0.65in}}
\toprule
\textbf{Type} & \textbf{GT} & \textbf{Matched} & \textbf{Missed} & \textbf{Recall} \\
\midrule
\PaperClipBreakdownRows
\bottomrule
\end{tabular}
}
\end{table}

The gradual-transition recall of 0.8089 is lower than the cut recall of
0.8991, but it is not the dominant source of the overall ClipShots gap.
The aggregate ClipShots precision is only 0.6498, indicating that open-world
effects and domain shift create many false-positive peaks.

\subsection{Training-Seed Robustness}
\label{sec:seed_robustness}

To quantify optimization variability, we retrained the BCE one-hot head
\PaperSeedCount{} times with training seeds 42, 43, and 44 while holding the
data seed, the training-video selection, and the sample cache identity fixed.
Each run selected its own post-processing configuration with the frozen
five-fold validation-only protocol of Section~\ref{sec:protocol} before any
test evaluation. Table~\ref{tab:seed_robustness} reports test means and
sample standard deviations. Seed-induced F1 variation is below 0.001 on SHOT
and BBC and below 0.005 on ClipShots, an order of magnitude smaller than the
video-sampling confidence intervals of Table~\ref{tab:confidence}; the
pipeline is therefore stable with respect to head-training randomness. The
absolute values differ from the controlled B4 numbers because this study
retrains the head from scratch and re-selects its protocol per seed (seed 42
selected $\sigma{=}1.0$ and $\theta{=}0.30$), so the two tables answer
different questions: Table~\ref{tab:ablation} isolates component
contributions on the historical cache, whereas this study isolates
training-seed sensitivity end to end.

\begin{table}[!t]
\caption{Three-Seed Robustness of the End-to-End Pipeline (Mean $\pm$ Sample Standard Deviation Over Training Seeds 42--44)}
\label{tab:seed_robustness}
\centering
\small
\renewcommand{\arraystretch}{1.08}
\resizebox{\columnwidth}{!}{%
\begin{tabular}{L{0.85in}C{1.0in}C{1.0in}C{1.0in}}
\toprule
\textbf{Dataset} & \textbf{F1} & \textbf{Precision} & \textbf{Recall} \\
\midrule
\PaperSeedRobustnessRows
\bottomrule
\end{tabular}
}
\end{table}

\subsection{Limitations}

The study has four important limitations. First, the headline numbers come
from the deployed checkpoint, whose result JSON files are tracked and
checksummed but whose test logits were not retained
(Section~\ref{sec:tracks}); the controlled track therefore carries every
logit-level claim, and the three-seed study of
Section~\ref{sec:seed_robustness} bounds optimization variability (F1 sample
standard deviation at or below 0.005 on every dataset) without reproducing
either checkpoint bit for bit. Second, the historical sample cache recorded
400 videos and 33,174 frames but not the exact selected IDs; strict
manifests fix this for future runs but cannot retroactively prove the old
cache contents. Third, validation calibration diagnostics reuse the
temperature-selection data. Fourth, the frozen representation does not fully
control cross-domain false positives on ClipShots.
